\documentclass[11pt]{amsart}
\usepackage[T1]{fontenc}
\usepackage{lmodern}
\usepackage{amsmath,amssymb,amsthm,mathtools}
\usepackage{xcolor}
\usepackage[colorlinks=true,linkcolor=blue!55!black,citecolor=blue!55!black,urlcolor=blue!55!black]{hyperref}

\newtheorem{theorem}{Theorem}[section]
\newtheorem{lemma}[theorem]{Lemma}
\newtheorem{proposition}[theorem]{Proposition}
\theoremstyle{definition}
\newtheorem{conjecture}[theorem]{Conjecture}
\newtheorem{definition}[theorem]{Definition}
\newtheorem{remark}[theorem]{Remark}

\title[Conjectural Binary-Kernel Rank]{Conjectural Infinite Binary-Kernel Rank for Three Subsequences of A182510\\
(Original Eight-Zero Problem Unresolved)}
\author{Xiaojian Zeng}
\address{PBC School of Finance, Tsinghua University, 100190}
\email{xiaojianzeng@foxmail.com}
\date{2026-09-04}
\hypersetup{pdftitle={Conjectural Infinite Binary-Kernel Rank for Three Subsequences of A182510 (Original Eight-Zero Problem Unresolved)},pdfauthor={Xiaojian Zeng}}

\begin{document}

\begin{abstract}
The original problem remains unresolved: it asks whether the OEIS A182510
recurrence has exactly eight zero terms.  The binary-kernel conjecture
formulated here is also unproved.  We study the integer
recurrence
\[
 a_0=0,\qquad a_1=1,\qquad
 a_n=(a_{n-1}\mathbin{\mathrm{XOR}}n)-a_{n-2},
\]
where XOR has infinite-width two's-complement semantics.  Exact reductions
connect its zero set to three unit-value subsequences.  Their canonical
binary-kernel evaluation matrices are nonsingular in every tested case
$3\leq d\leq11$, reaching dimension $2047$, and independent exact
implementations reproduce the certificates.  Separate adversarial work finds
large forced powers of two in these determinants and invalidates two naive
proof mechanisms without producing a dependence.  This evidence motivates
Conjecture~\ref{conj:C002}, which predicts nonsingularity at every depth.  It
does not imply the original eight-zero assertion.  No case $d\geq12$ has been
tested, no determinant recurrence is known, and neither finite computation
nor the recorded bounded literature search supplies a proof.
\end{abstract}

\maketitle

\section{Original problem and verified frontier}

For $x\in\mathbb Z$ and $n\in\mathbb N$, infinite-width two's-complement
XOR is fixed by
\[
x\mathbin{\mathrm{XOR}}n=
\begin{cases}
x\mathbin{\mathrm{xor}}n,&x\geq0,\\
-1-\bigl((-1-x)\mathbin{\mathrm{xor}}n\bigr),&x<0,
\end{cases}
\]
where the operation on the right is ordinary finite binary XOR.  There is no
word-size truncation or unsigned coercion.  The original problem asks whether
\begin{equation}\label{eq:original-target}
 \{n\in\mathbb N:a_n=0\}=\{0,6,11,23,36,42,48,54\}.
\end{equation}
The live OEIS entry gives the same recurrence and labels the assertion that
the sequence has eight zeros as a conjecture; its linked table covers indices
$0$ through $10000$ \cite{OEISA182510}.  A targeted search on 2026-09-04
found no accepted proof or counterexample.  This is a bounded open-status
check, not an exhaustive survey: the OEIS entry is a curated database record,
not a peer-reviewed proof; the indexed web searches performed here cannot be
exhaustive; and a supplied
commit-pinned formalization could not be inspected.
Accordingly, \eqref{eq:original-target} remains unresolved in this note.

Two exact finite frontiers are useful but do not settle the problem.  First,
independent recurrence implementations find precisely the eight displayed
zeros through $n=5\cdot10^{10}$.  Second, at dyadic level $33$, any integer
zero must occupy one of $354$ completely enumerated residue classes modulo
$3\cdot2^{33}$; the least additional representative is $123{,}064{,}307$.
Both statements concern finite, reproducible computations and make no claim
about the infinite tail.

The closest literature located for the method is the theory of
$k$-regular sequences.  It characterizes $k$-regularity by finite generation
of the module spanned by the $k$-kernel
\cite{AlloucheShallit1992,AlloucheShallit2003}; the 1992 paper also uses
nonvanishing finite evaluation determinants for a different sequence.  Work
on compatible $2$-adic dynamics supplies broader context for arithmetic and
bitwise maps \cite{Anashin2002}.  None of the inspected sources states the
recurrence-specific conjecture below or classifies its three derived
subsequences.

\section{Research routes and failure diagnosis}

The investigation kept proof and counterexample routes separate.  The main
outcomes, including failures, are as follows.

\begin{enumerate}
\item A signed-size route gives
\(
 |(x\mathbin{\mathrm{XOR}}n)-x|\leq n
\)
and hence $|a_n|\leq n(n+1)/2$.  It cannot exclude zeros: the homogeneous
recurrence has a six-periodic impulse response, so absolute-value
accumulation discards the sign and low-bit data needed for a lower bound.

\item A modular route succeeds in proving a universal period modulo every
power of two (Section~\ref{sec:partial}), but it leaves a binary forest of
compatible zero residues.  At level $33$ there are $346$ non-known residue
classes in addition to the eight known zeros.  No height function, finite
automaton, or uniform escape argument classifies their infinite lifts.

\item A direct counterexample search uses two different exact update formulas
and guarded signed arithmetic.  The implementations agree through
$5\cdot10^{10}$ and find no ninth zero.  The stopping point is a computational
budget boundary, not evidence for a universal tail statement.

\item Three fixed-window lift predictors were attacked on parent levels
$1$ through $24$, for every window width $1$ through $12$.  All $36$
feature/width combinations contain exact collisions: identical retained
features can lead to different child masks.  This rules out those particular
static representations, but not a level-dependent state or a state carrying
the propagated borrow.

\item Exact three-step and bit-stripping identities reduce zeros to unit
values in three residue-class subsequences.  Low-dimensional linear
binary-kernel closure then fails: each of the three relevant rational
$2$-kernels has dimension at least $2047$.  This finite obstruction suggests
the conjecture in Section~\ref{sec:conjecture}, but it neither gives infinite
rank nor classifies unit hits.

\item A second stripping isolates an internal subtraction and its borrow.
For every fixed low-bit window there are same-window, same-phase integer
states requiring different corrections.  This eliminates an all-state static
affine ansatz; it does not eliminate a bit-serial transducer or a closure
specific to the distinguished orbit.

\item Binary-parent column differences prove substantial all-depth
$2$-adic divisibility of the conjectured determinants.  Exact prime factors
present at one depth can disappear at the next, refuting a naive recursion
that retains the preceding determinant as a factor.  Divisibility by a large
power of two is compatible with both a zero and a nonzero determinant.
\end{enumerate}

The smallest explicit unresolved kernel for the original problem is the
following lift-tree classification.  Let
\[
 Z_k=\{0\leq r<3\cdot2^k:a_r\equiv0\pmod{2^k}\}.
\]
Each node has only the possible children $r$ and $r+3\cdot2^k$.
Ordinary integer zeros are exactly the infinite paths whose representatives
eventually remain constant.  It remains to exclude every eventually constant
path except those stabilizing at the eight indices in
\eqref{eq:original-target}.  This finite-at-each-level reformulation is more
testable than an undirected tail search, but it is still equivalent to the
missing universal assertion rather than a solution.

\section{Verified partial results}\label{sec:partial}

This section records only identities with complete elementary derivations or
exact finite certificates.  None resolves the original problem or the
conjecture in Section~\ref{sec:conjecture}.

\subsection{Universal modular period and the zero-lift tree}

Put $M=2^k$ and, on $(\mathbb Z/M\mathbb Z)^2$, define
\[
 T_j(x,y)=\bigl(y,(y\mathbin{\mathrm{XOR}}j)-x\bigr),\qquad
 P_M=T_{M-1}\cdots T_0.
\]
A high-bit translation induction gives
\begin{equation}\label{eq:period}
 P_M^3=\operatorname{id},\qquad
 a_{n+3\cdot2^k}\equiv a_n\pmod{2^k}
 \quad(k,n\geq0).
\end{equation}
For completeness, at the lift from $M$ to $2M$ let
$H_{(\alpha,\beta)}(x,y)=(x+\alpha M,y+\beta M)$,
$e=(0,1)$, and $A(\alpha,\beta)=(\beta,\alpha+\beta)$ over
$\mathbb F_2$.  Direct separation of the lower and upper bits gives
\[
 T_{j+M}=H_eT_j,\qquad T_jH_v=H_{Av}T_j,\qquad
 A^3=I,\qquad I+A+A^2=0.
\]
If $Q=T_{M-1}\cdots T_0$, commuting the translations through the upper
half of the word yields $P_{2M}=H_sQ^2$ with
$s=\sum_{i=0}^{M-1}A^ie$.  Cubing, the translation part vanishes by
$I+A+A^2=0$.  The induction hypothesis makes $Q^3$ a translation on each
four-point fiber over modulus $M$, and every such translation has square
equal to the identity.  Thus $Q^6=I$ and the induction closes, including
$M=1$.  Cyclically rotated time words are conjugate, which gives the second
identity in \eqref{eq:period} for every starting phase.

Reduction of \eqref{eq:period} from level $k+1$ to $k$ gives the parent map
for $Z_k$.  A path eventually fixed at $n$ lies in $Z_k$ for arbitrarily
large $k$, so $2^k\mid a_n$ for arbitrarily large $k$ and hence $a_n=0$.
The converse is immediate.  This proves the lift-tree equivalence used above,
but supplies no classification of non-known paths.

\subsection{Local zero tests and unit-value coordinates}

XOR is an involution, so for every $n\geq2$,
\begin{equation}\label{eq:adjacent-xor}
 a_n=0
 \quad\Longleftrightarrow\quad
 a_{n-2}\mathbin{\mathrm{XOR}}a_{n-1}=n.
\end{equation}
Consequently the two preceding values have the same sign, agree above the
highest set bit of $n$, and differ by at most $n$.  Modulo two, the parity
word is $(0,1,1,1,0,0)^\infty$, so a zero index is congruent to $0$, $4$,
or $5$ modulo $6$.

Eliminating adjacent recurrence terms and using the last binary carry gives
the exact equivalent unit-value classifications
\begin{align}
 \{q\geq0:a_{6q+3}=-1\}&=\{0,1,6,7,8,9\},\label{eq:unit-w}\\
 \{q\geq0:a_{6q+1}=1\}&=\{0\},\label{eq:unit-u}\\
 \{q\geq0:a_{6q+2}=1\}&=\{1,3\}.
 \label{eq:unit-v}
\end{align}
Their conjunction is equivalent to \eqref{eq:original-target}; it is not a
proof of any of the three equalities.  The proved period modulo four makes
the following integer-valued coordinates natural:
\begin{equation}\label{eq:uvw}
 U_q=\frac{a_{6q+1}-1}{4},\qquad
 V_q=\frac{a_{12q+8}-1}{4},\qquad
 W_q=\frac{a_{6q+3}+1}{4}.
\end{equation}
These isolate the nontrivial unit hits in
\eqref{eq:unit-w}--\eqref{eq:unit-v} and expose their binary descendants.

\subsection{Finite kernel-rank certificates and forced divisibility}

For $S\in\{U,V,W\}$, evaluating the first $2^d-1$ standard binary-kernel
subsequences at $j=0,\ldots,2^d-2$ gives the matrix $K_d^S$ used below.
At $d=11$, exact modular elimination gives
\begin{align*}
 \det K_{11}^U&\equiv329819\pmod{1000003},\\
 \det K_{11}^V&\equiv646456\pmod{1000003},\\
 \det K_{11}^W&\equiv520101\pmod{1000003}.
\end{align*}
Each residue is nonzero.  Therefore the corresponding integer determinant is
nonzero, and the first $2047$ binary-kernel subsequences of each sequence are
linearly independent over $\mathbb Q$.  This proves only the finite lower
bound $2047$.

There is also an all-depth obstruction to a simple odd-determinant argument.
Writing $D_E(S)=\det K_{E+1}^S$, binary-parent column differences followed by
a parity-rank estimate give, for every $E\geq2$,
\begin{align}
 2^{(4E-2)2^E-3E+1}&\mid D_E(V),\label{eq:div-v}\\
 2^{(4E-4)2^E-3E+3}&\mid D_E(U),D_E(W).
 \label{eq:div-uw}
\end{align}
The normalized parity after the column differences depends only on row depth
and column scale.  Factoring the forced powers of two and applying the
elementary fact that an $m$-square integer matrix of rank at most $\rho$
modulo two has determinant divisible by $2^{m-\rho}$ gives
\eqref{eq:div-v}--\eqref{eq:div-uw}.  Since zero is divisible by every power
of two, these formulas do not imply nonsingularity.

\section{Robust candidate conjecture}\label{sec:conjecture}

The following item passed a fresh independent conjecture audit as a
precise, reproducible, adversarially tested candidate.  It remains unproved.
Its mathematical scope is stated precisely below.

\begin{conjecture}[Registry C002]\label{conj:C002}
For every $q\in\mathbb{N}$, define
\begin{align*}
 U_q&=(a_{6q+1}-1)/4,&
 V_q&=(a_{12q+8}-1)/4,&
 W_q&=(a_{6q+3}+1)/4.
\end{align*}
For every $S\in\{U,V,W\}$ and every integer $d\ge 3$, let $K_d^S$ be
the $(2^d-1)\times(2^d-1)$ integer matrix whose rows are indexed, in
lexicographic level order, by pairs $(e,r)$ satisfying
\[
 0\le e<d,
 \qquad 0\le r<2^e,
\]
whose columns are indexed by integers $j$ satisfying $0\le j<2^d-1$, and
whose entries are
\[
 K_d^S[(e,r),j]=S_{2^e j+r}.
\]
Then $\det K_d^S\ne0$.
\par\smallskip
\noindent\emph{Canonical source formulation.}
{\footnotesize\ttfamily\sloppy
\detokenize{For every q\in\mathbb{N}, define U_q=(a_{6q+1}-1)/4, V_q=(a_{12q+8}-1)/4, and W_q=(a_{6q+3}+1)/4. For every S\in\{U,V,W\} and every integer d\ge 3, let K_d^S be the (2^d-1)\times(2^d-1) integer matrix with rows indexed, in lexicographic level order, by pairs (e,r) satisfying 0\le e<d and 0\le r<2^e, columns indexed by integers j satisfying 0\le j<2^d-1, and entries K_d^S[(e,r),j]=S_{2^e j+r}. Then \det K_d^S\ne0.}\par}
\par\smallskip
\noindent\emph{Plain-language statement.} For every nonnegative integer $q$,
define
\begin{align*}
 U(q)&=(a(6q+1)-1)/4,\\
 V(q)&=(a(12q+8)-1)/4,\quad\text{and}\\
 W(q)&=(a(6q+3)+1)/4.
\end{align*}
Choose any $S$ among $U$, $V$, and $W$ and any integer $d$ at least $3$.
Form the square integer matrix of dimension $2^d-1$ whose rows are indexed
by all pairs $(e,r)$ with $0\le e<d$ and $0\le r<2^e$, whose columns are
$j=0,\ldots,2^d-2$, and whose $((e,r),j)$ entry is $S(2^e j+r)$.  The
conjecture says that this matrix has nonzero determinant for every such $S$
and $d$.
\end{conjecture}

The conjecture is nontrivial already at its boundary: all three $d=1$
matrices are singular, while the $d=2$ matrices for $V$ and $W$ are
singular.  Thus $d\geq3$ is the smallest uniform threshold for the three
sequences.  It is also not a reformulation of the original problem.
Conjecture~\ref{conj:C002} does not classify the zeros of $U,V,W$ and does
not imply \eqref{eq:original-target}.

If the conjecture is true, the rational $2$-kernel of each of $U,V,W$ has
infinite dimension, so none is $2$-regular.  It would rule out every
finite-dimensional linear binary-kernel description of these unit-hit
coordinates and redirect the original structural search toward nonlinear or
unbounded carry data.  If the conjecture is false, a singular canonical minor
would not by itself imply finite kernel rank: an exact null relation should be
tested on additional columns to distinguish a whole-subsequence dependence
from failure of this chosen family of minors.  Alternative square minors
could still witness unbounded rank.

\section{Exact evidence and reproducibility}

All decisive computations use integer or finite-field arithmetic; no
floating-point inference is used.  For $3\leq d\leq10$, every matrix for
$S\in\{U,V,W\}$ has a nonzero determinant modulo both $1000003$ and
$1000033$.  For $d=11$, all three determinants are nonzero modulo $1000003$.
The corresponding dimensions are $2^d-1$, culminating at $2047$.  The
smallest claimed matrices were also computed over the integers:
\begin{align*}
 \det K_3^U&=4029874176,\\
 \det K_3^V&=-146688001966080,\\
 \det K_3^W&=73891840.
\end{align*}

The producer generates the recurrence using direct signed XOR and performs
unnormalized modular elimination.  The independent checker regenerates every
sequence value from
\[
 x\mathbin{\mathrm{XOR}}n=x+n-2(x\mathbin{\mathrm{AND}}n)
\]
and uses a separately normalized elimination.  It reconstructs every matrix
and recomputes every determinant residue.  The largest recurrence frontier is
$25{,}153{,}532$; a conservative operation bound remains below the signed
$64$-bit limit.  A third, arbitrary-precision implementation with exact
domain determinants checks the boundary and small claimed levels.

The determinant evidence, adversarial normalized-layer evidence, and their
independent checkers are stored in
\begin{center}
\begin{minipage}{0.94\textwidth}
\small
\path{experiments/kernel-rank-evidence.json},
\path{experiments/kernel-rank-d11-evidence.json},
\path{experiments/handoff-extensions-evidence.json},
\path{experiments/verify_kernel_rank.py}, and
\path{experiments/verify_handoff_extensions.py}.
\end{minipage}
\end{center}
From the project root, the principal replay commands are
\begin{verbatim}
python3 experiments/verify_kernel_rank.py \
  experiments/kernel-rank-evidence.json
python3 experiments/verify_kernel_rank.py \
  experiments/kernel-rank-d11-evidence.json
python3 experiments/verify_handoff_extensions.py \
  experiments/handoff-extensions-evidence.json
python3 experiments/verify_structural_hints.py \
  experiments/structural-hints-evidence.json
\end{verbatim}
The saved kernel-rank environment records Python $3.10.12$.  It also records
NumPy $2.2.6$.  The handoff checker additionally uses exact SymPy domain
determinants.  Successful replay certifies only the serialized finite scope.

For context, the separate original-problem computations cover the direct
sequence through $5\cdot10^{10}$, the full distinguished modular orbit through
$k=33$, all initial states through $k=8$, all phases through $k=5$, and
$748{,}980$ exhaustive plus $100{,}000$ seeded high-bit lift tests.  They are
exact within their declared finite ranges but are not evidence for an
unbounded quantifier.

\section{Adversarial falsification and literature boundary}

The dedicated adversarial route differs from merely extending the positive
rank table.  It included the following attacks.

\begin{itemize}
\item The excluded levels $d=1,2$ were evaluated.  Their singularities show
that the threshold in Conjecture~\ref{conj:C002} cannot be lowered uniformly.

\item Searches tested duplicate or linearly dependent low-level rows and
forced congruence degeneracy.  They also tested exact rational dependence in
every claimed finite case.  No dependence was found.

\item Binary-parent differences were normalized through depth $E=8$, for
$781{,}828$ exact entries.  Exact integer determinants through $E=4$ were
computed by both fraction-free Bareiss elimination and an independent SymPy
implementation.  These tests confirm the forced divisibility
\eqref{eq:div-v}--\eqref{eq:div-uw} without converting it into
nonsingularity.

\item Prime factors of $D_2(U)$, $D_2(V)$, and $D_2(W)$ disappear from the
corresponding $D_3$ in explicit examples.  This refutes a naive nested
determinant factorization.  The propagated-borrow analysis separately refutes
an all-state affine correction based on any fixed low-bit window.

\item Optimization extremizers and random input search are inapplicable:
each pair $(S,d)$ determines a single integer matrix.  Multiple modular primes
are useful exact prefilters, not randomized counterexamples.
\end{itemize}

The fastest future falsifier is concrete.  Generate a specified $K_d^S$,
use several modular determinants as inexpensive singularity prefilters, then
reconstruct a nonzero rational or integer null vector and independently check
that its product with the serialized integer matrix is exactly zero.  A zero
determinant modulo one prime, or even several primes, is not a refutation.

The novelty search used the exact OEIS identifier and recurrence, the terms
``$2$-kernel'', ``$2$-regular'', determinant, nonsingularity, canonical
evaluation minor, and the patterns $a(6n+1)$ and $a(12n+8)$.  It also checked
the primary Allouche--Shallit papers and their determinant precedent
\cite{AlloucheShallit1992,AlloucheShallit2003}.  No equivalent of the fully
specified matrix statement was found in that recorded search as of
2026-09-04.  This wording is deliberately bounded: incomplete indexing and
the inaccessible pinned formalization preclude an absolute priority or
novelty claim.

\section{Limitations, unsearched scope, and research agenda}

Conjecture~\ref{conj:C002} is unproved, and the original eight-zero problem
remains unresolved.  The finite determinant data stop at $d=11$; every
$d\geq12$ is untested.  No exact determinant recurrence, higher normalized
quotient-layer argument, or universal odd residual has been obtained.  Exact
null-vector searches after a multi-prime modular singularity, alternative
evaluation minors, and tests for whole-subsequence dependence also remain to
be done.  Nonlinear finite-carry automata and level-dependent state
descriptions lie outside the conjecture.

For the original problem, no argument excludes a zero beyond
$5\cdot10^{10}$, and no theorem classifies all compatible zero-residue paths.
The longest observed non-known run of lower lift edges through level $33$ is
$11$, but no uniform bound follows.  The most direct structural next step is
to propagate the exact borrow through another bit-stripping layer on the
distinguished orbit and test whether a conflict-free bit-serial state closes.
For Conjecture~\ref{conj:C002}, the corresponding next step is to derive or
falsify a recurrence for the next normalized determinant quotient.

The literature review is current only to the recorded 2026-09-04 search and
does not guarantee exhaustive novelty.  The computational certificates test
large but finite objects.  Neither those certificates, the independent
conjecture audit, nor the robust-candidate designation is a probability or a
truth certificate.  No proof assistant was used.

\section*{Data, code, and computational disclosure}

The project dossier contains the frozen problem statement, route registry,
gap ledger, failure report, exact JSON evidence, producer programs, and
independent checkers.  The large direct search uses guarded signed arithmetic;
the determinant work uses integer and prime-field arithmetic; and the
independent implementations use a distinct XOR identity and elimination
normalization.

AI systems assisted with proposing and attacking research routes, writing and
checking computational code, organizing derivations, performing a fresh
conjecture audit in a context separate from the owner, and drafting this note.
Every reported mathematical claim is tied to an explicit derivation or exact
reproducible artifact; finite computation is not treated as proof of an
infinite statement.  The author remains responsible for the mathematical
content.  No proof assistant was used.

\bibliographystyle{amsplain}
\bibliography{references}

\end{document}
